\begin{solapartat}
El diagrama de forces serà el següent:

\noindent\begin{minipage}[c]{0.55\linewidth}
\figura[0.9]{sol_fig1.png}
\end{minipage}\hfill
\begin{minipage}[c]{0.3\linewidth}
\figura[0.8]{sol_fig2.png}
\end{minipage}\hfill
\begin{minipage}[c]{0.1\linewidth}
\punts{0,3}
\end{minipage}

Per fer el dibuix cal tenir en compte que la càrrega del electró és negativa. $\vec{F}_E$, és la força deguda al camp elèctric i $\vec{F}_B$ és la deguda al camp magnètic.

Per tal que els electrons no es desviïn al travessar aquesta regió les dues forces han de ser iguals:
\[ \left.\begin{aligned} F_E &= q_e\,E\\ F_B &= v\,q_e\,B \end{aligned}\right\} \Rightarrow q_e\,E = v\,q_e\,B\ \text{\punts{0,3}} \Rightarrow v = \frac{E}{B}\ \text{\punts{0,2}} = \frac{250}{0,04} = 6,25\cdot 10^{3}\ \mathrm{m/s}\ \text{\punts{0,2}} \]
\end{solapartat}

\begin{solapartat}
Tal com es pot veure del gràfic anterior la força deguda al camp magnètic es paral·lela al pla XY, per tan l'electró farà un moviment circular en un pla paral·lel al XY \punts{0,2} i en el sentit horari. \punts{0,1}

La força $\vec{F}_B$ és la que proporcionarà l'acceleració centrípeta per fer girar l'electró, per trobar el radi de gir tindrem:
\[ m\,\frac{v^2}{r} = v\,q_e\,B \Rightarrow r = \frac{mv}{q_eB} = 1,78\cdot 10^{-6}\ \mathrm{m}\ \text{\punts{0,2}} \]

La freqüència angular de rotació és:
\[ \omega = \frac{v}{r} = \frac{q_eB}{m}\ \text{\punts{0,3}} \Rightarrow \nu = \frac{\omega}{2\pi} = \frac{q_eB}{2\pi m} = 1,12\cdot 10^{9}\ \mathrm{Hz}\ \text{\punts{0,2}} \]
\end{solapartat}
